% !TEX TS-program = lualatex
Программное обеспечение сегодня, от простых утилит до сложных многопользовательских сервисов, в значительной степени основывается на моделях клиент–серверной архитектуры. Этот подход позволяет задействовать многочисленные вычислительные узлы, географически распределённые и объединённые в единую систему, что делает реальным выполнение задач, ранее невозможных или неэффективных в рамках монолитной архитектуры. Существуют различные причины использования такого подхода к развертыванию приложений, например:

\begin{enumerate}
    \item Необходимость в высоких вычислительных мощностях. Многие современные приложения, особенно в областях обработки больших данных, машинного обучения, видеостриминга или онлайн-игр, требуют значительных ресурсов. Такой подход позволяет гибко масштабировать инфраструктуру, автоматически увеличивая или уменьшая объём задействованных ресурсов в зависимости от текущей нагрузки, что делает приложения более производительными и надёжными.
    \item Требования к безопасности и защите коммерческой тайны диктуют необходимость особого подхода к хранению и обработке данных. Распределённая архитектура позволяет изолировать критически важные компоненты, использовать многоуровневые механизмы защиты, шифрования и контроля доступа. Кроме того, такие системы легче адаптируются под нормативные требования различных юрисдикций (например, законы о хранении персональных данных).
    \item В распределённых системах можно обновлять отдельные модули или сервисы без остановки всей системы, что повышает устойчивость и доступность приложения. Это особенно важно в условиях непрерывной доставки и интеграции, ставших стандартом в современной разработке.
\end{enumerate}

В рамках настоящей работы была выбрана архитектура, ориентированная на модель клиент–сервер. Клиенту, как конечному пользователю, не требуется взаимодействовать с компонентами программного обеспечения, предназначенными для обучения искусственного интеллекта. Модели, поставляемые в составе программного решения, могут обновляться независимо и всегда поддерживаются в актуальном состоянии. Кроме того, использование браузера устраняет необходимость установки приложения на устройство пользователя и потенциально расширяет перечень поддерживаемых платформ для игры.
